% TOP_PROBLEM: {"problemId": "problem.coding-theory-and-combinatorics-problem-cc-6-existence-of-finite-projective-planes-of-non-prime-power-orders", "canonicalTitle": "Coding Theory and Combinatorics Problem CC-6: Existence of Finite Projective Planes of Non-Prime-Power Orders", "releaseRank": 404, "primaryDomain": "cryptography_coding_information_optimization", "primaryDomainLabel": "Cryptography, coding, information & optimization", "displayStatus": "open", "releaseStatus": "open"}
% !TeX program = lualatex
% ATTEMPT_STATUS: unresolved
\documentclass[11pt]{article}
\usepackage[margin=25mm]{geometry}
\usepackage{fontspec}
\setmainfont{Latin Modern Roman}
\usepackage{amsmath,amssymb,mathtools}
\usepackage{unicode-math}
\setmathfont{Latin Modern Math}
\usepackage{xurl,hyperref}
\hypersetup{colorlinks=true,urlcolor=blue}
\setcounter{secnumdepth}{0}
\setlength{\parindent}{0pt}
\setlength{\parskip}{5pt}
\setlength{\emergencystretch}{3em}
\begin{document}
\section*{\texorpdfstring{Coding Theory and Combinatorics Problem CC-6: Existence of Finite Projective Planes of Non-Prime-Power Orders}{Coding Theory and Combinatorics Problem CC-6: Existence of Finite Projective Planes of Non-Prime-Power Orders}}\label{404}
\subsection{Definitions and mathematical statement}
A finite projective plane is a finite point--line incidence structure in which two distinct points lie on a unique line, two distinct lines meet at a unique point, and four points exist with no three collinear. Its order $q$ means each line contains $q+1$ points; it then has $q^2+q+1$ points and as many lines. The problem asks whether a plane exists of order $12$, or more generally of any order $q\ge2$ which is not a prime power and is not excluded by the Bruck--Ryser--Chowla obstruction. In particular, for $q\equiv1,2\pmod4$ that obstruction requires $q$ to be a sum of two integer squares. Satisfying the obstruction is only a necessary condition, not a construction. The order-12 question is a particular case of the broader existence question.
\par\smallskip\textit{Formulation and concept references:} \href{https://arxiv.org/pdf/0811.1254}{[S1]}.

\subsection{Short English statement}
Can a finite projective plane have order twelve, or any other allowable order that is not a prime power?
\subsection{Research attempt}
\paragraph{Scope, process, and first gap.}
This GPT-6 Astra Ultra attempt, dated 27 September 2026, studies the incidence-matrix and Latin-square formulations, then tests replacing a finite field by arithmetic modulo twelve. It derives equivalent finite construction targets and checks necessary conditions. It does not construct or exclude a plane of order twelve.

The original question includes arbitrary projective planes, not just planes arising from fields or satisfying Desargues's theorem. Thus proving that finite fields have prime-power cardinality is not a nonexistence proof for the requested planes. The first gap is constructing, or excluding, an unrestricted incidence structure at a permitted non-prime-power order.

\paragraph{The incidence matrix and its converse.}
Put $v=q^2+q+1$, and let $A$ be the $v$ by $v$ binary matrix whose rows are lines and columns are points. The plane axioms imply
\[
 AA^{\mathsf T}=qI+J,
 \tag{404.1}
\]
where $J$ is the all-ones matrix. The diagonal counts $q+1$ points on a line; each off-diagonal entry counts the unique intersection of two lines.

Conversely, suppose a binary square matrix of this size satisfies (404.1). Its row sums are $q+1$, because a binary row has squared norm equal to its sum. The matrix $qI+J$ is positive definite, so $A$ is invertible over the reals. Moreover,
\[
 A\bigl(A^{\mathsf T}\mathbf1-(q+1)\mathbf1\bigr)
 =(q+v)\mathbf1-(q+1)^2\mathbf1=0.
\]
Invertibility gives $A^{\mathsf T}\mathbf1=(q+1)\mathbf1$: column sums also equal $q+1$. Both $AJ$ and $JA$ therefore equal $(q+1)J$. Conjugating (404.1) by $A$ now gives
\[
 A^{\mathsf T}A=A^{-1}(qI+J)A=qI+J.
\]
Thus every pair of distinct points lies on exactly one line.

The nondegeneracy axiom also follows for $q\ge2$. Choose two points on a line and a point outside that line; these form a triangle. Its three sides have union of size $3(q+1)-3=3q$. There are
\[
 v-3q=(q-1)^2>0
\]
points outside those sides. Any such point completes a quadrangle. Hence the binary Gram equation is an exact formulation, not merely a relaxation.

For $q=12$ this asks for a $157$ by $157$ binary matrix, each row of weight $13$, with pairwise row intersections exactly one. Positivity of the real matrix on the right alone does not construct binary factors.

\paragraph{What the simplest determinant test misses.}
The eigenvalues of $qI+J$ are $q$ with multiplicity $v-1$ and $q+v=(q+1)^2$ once. Therefore
\[
 (\det A)^2=q^{\,v-1}(q+1)^2.
 \tag{404.2}
\]
Since $v-1=q(q+1)$ is always even, the right side is an integer square for every integer $q$. The elementary determinant-square test therefore excludes no order here.

The stronger Bruck--Ryser--Chowla condition in [S1] uses more than (404.2). For $q\equiv1,2\pmod4$, it requires a representation as a sum of two integer squares. Order six fails that condition; order ten passes because $10=1^2+3^2$, yet no order-ten plane exists, as established by the work discussed in Lam's author account [E1]. This proves that passing the condition is not sufficient in general.

Order twelve is congruent to zero modulo four, so the displayed sum-of-two-squares restriction does not apply to it. The fact that twelve is not a sum of two squares is therefore irrelevant to this particular obstruction. No stronger exclusion of twelve is obtained from the calculations above.

\paragraph{What fields construct, and what a ring does not.}
For a finite field $\mathbb F_q$, take one-dimensional subspaces of $\mathbb F_q^3$ as points and two-dimensional subspaces as lines, with containment as incidence. Two distinct one-dimensional subspaces span a unique two-dimensional one. Two distinct two-dimensional subspaces intersect in dimension one. There are $(q^3-1)/(q-1)=q^2+q+1$ points and $q+1$ points on each line; duality gives the line count. The coordinate points together with the span of $(1,1,1)$ give a quadrangle.

A finite field has prime characteristic $p$ and is a finite-dimensional vector space over $\mathbb F_p$, so its cardinality is $p^r$. This accounts for the orders supplied by this construction. It says nothing about whether another construction gives a non-Desarguesian plane at a non-prime-power order.

The naive replacement $\mathbb F_q\rightsquigarrow\mathbb Z/12\mathbb Z$ fails concretely. Use primitive triples modulo multiplication by units as proposed projective points and primitive linear forms as proposed lines. The two distinct points
\[
 P=(1,0,0),\qquad Q=(1,6,0)
\]
lie on both proposed lines
\[
 L_1:\ z=0,\qquad L_2:\ 2y+z=0.
\]
All four coordinate triples are primitive. The points cannot be unit multiples because their first coordinate is one and their second coordinates differ. The line forms are not unit multiples either; their last coordinate forces the multiplier to be one. The two line sets differ, for example at $(0,1,0)$. Nevertheless both displayed points satisfy both equations modulo twelve.

Thus even this restricted ring construction violates unique joins. The zero divisor $6$ is responsible for $2\cdot6=0$. This defeats the proposed construction, not every possible plane of order twelve.

\paragraph{An equivalent target using Latin squares.}
Delete a line at infinity from a plane of order $q$. The remaining $q^2$ points form an affine plane. The other lines fall into $q+1$ parallel classes, each containing $q$ disjoint lines of size $q$. Lines from different classes meet once.

Choose two classes as rows and columns. Each affine point is identified by its row and column. For each remaining class, label its $q$ lines with $q$ symbols and put the label of the line through a cell in that cell. Every symbol occurs once in each row and column, so this is a Latin square. Two such squares are orthogonal because each line of one class meets each line of the other once. We obtain $q-1$ mutually orthogonal Latin squares.

Conversely, from $q-1$ mutually orthogonal Latin squares of order $q$, use cells as points and take rows, columns, and symbol classes as lines. These give $q+1$ parallel classes. Distinct classes have lines meeting exactly once, by the Latin and orthogonality conditions.

For a fixed cell, its $q+1$ incident lines have no other intersections with one another. Their remaining points number
\[
 (q+1)(q-1)=q^2-1,
\]
so they cover every other cell exactly once. This proves the unique line axiom for all affine point pairs, rather than assuming it. Add one point at infinity for each parallel class and one line containing all those new points. The projective axioms follow.

Consequently order twelve is equivalent to a complete set of eleven mutually orthogonal Latin squares of order twelve. Producing one Latin square or a single orthogonal pair is far short of this exact target; the number of simultaneously compatible squares is essential.

\paragraph{Verification and outcome.}
Finite checks construct the field planes for $q=2,3,5$, verify both Gram identities, and test the corresponding complete Latin-square families. They verify the modulo-twelve counterexample with exact modular arithmetic. These are sanity checks of the constructions and equivalences, not a search through the order-twelve incidence matrices.

\textbf{UNRESOLVED.} The binary matrix and complete Latin-square targets are exact, and the determinant and ring shortcuts have been audited. An unrestricted order-twelve matrix satisfying (404.1), or an obstruction applying to every such matrix, remains missing. The broader allowable non-prime-power question is likewise unanswered.

\subsection{Sources}
\begin{itemize}
\item[S1]Huber, Coding theory and algebraic combinatorics (2008), Remark 1.1.
\url{https://arxiv.org/pdf/0811.1254}.
\textit{Formulation source.}
\item[E1]Clement W. H. Lam, The Search for a Finite Projective Plane of Order 10, American Mathematical Monthly 98 (1991), 305--318, author-hosted account.
\url{https://www.cecm.sfu.ca/organics/papers/lam/index.html}.
\textit{Primary account of the established order-ten nonexistence result; it does not settle order twelve.}
\end{itemize}
\end{document}
